\section{Diseño del Sistema y Arquitectura}
\label{sec:arquitectura}

\subsection{Modelo C4}

La arquitectura se documenta en tres niveles del modelo C4 (\texttt{docs/arquitectura/README.md}):

\begin{itemize}[nosep]
    \item \textbf{Nivel 1 (Contexto):} el sistema se relaciona con cuatro tipos de actores (Estudiante,
    Docente/Jurado, Coordinador, Administrador) y un sistema externo (API de universidades de Hipo Labs,
    consumida vía \texttt{ExternalApiServiceImpl} con caché Redis de 10 minutos).
    \item \textbf{Nivel 2 (Contenedores):} cuatro contenedores desplegables --- Frontend Angular (SPA),
    Backend Spring Boot (API REST), PostgreSQL 15 (persistencia), Redis 7 (caché), coordinados por nginx
    como proxy inverso.
    \item \textbf{Nivel 3 (Componentes):} 25 controladores REST, organizados por dominio (Auth,
    Solicitudes, Jurados, Cronograma, Evaluaciones, Actas, Notificaciones, Reportes, Salas, Usuarios,
    Anteproyectos, y otros de soporte), cada uno delegando a una capa de servicio y, para la lógica
    académica compleja, a procedimientos almacenados PL/pgSQL.
\end{itemize}

\subsection{Resumen de los siete Registros de Decisión Arquitectónica (ADR)}

\begin{table}[H]
\centering
\small
\begin{tabular}{p{1.6cm}p{3.2cm}p{7cm}}
\toprule
\textbf{ADR} & \textbf{Decisión} & \textbf{Justificación resumida} \\
\midrule
ADR-001 & Arquitectura general de 3 capas & Separación Angular/Spring Boot/PostgreSQL para independencia de despliegue y escalado \\
ADR-002 & Autenticación JWT + cookies HTTP-Only & Stateless para escalar horizontalmente sin sesión compartida; cookie HTTP-Only mitiga XSS sobre el token \\
ADR-003 & Frontend Angular (SPA) & Reactividad y componentización para un dominio con múltiples roles y vistas condicionales \\
ADR-004 & Seguridad alineada a OWASP & JWT de 7 claims, rate limiting, cabeceras HSTS/CSP/X-Frame-Options, BCrypt \\
ADR-005 & Control de acceso por permisos dinámicos (BD), no roles estáticos en código & Reemplaza el RBAC fijo con \texttt{@PreAuthorize} descrito originalmente en ADR-004 por una verificación contra \texttt{permisos}/\texttt{rol\_permisos}, para que cambiar qué rol puede qué requiera solo datos, no redespliegue \\
ADR-006 & Estrategia híbrida JPA + procedimientos almacenados (separación CRUD/SP) & CRUD simple vía Spring Data JPA; lógica académica compleja (notas ponderadas, reportes, firmas) vía PL/pgSQL para evitar SQL dinámico. Ampliada en la Fase 8 con la decisión explícita de mantener la convención Flyway en vez de espejar en \texttt{db/procs/} \\
ADR-007 & Despliegue con Docker Compose, imágenes ancladas por digest SHA-256 & Reproducibilidad determinista. Ampliada en la Fase 8 con la decisión de Railway como proveedor de producción (Sección~\ref{sec:implementacion}) \\
\bottomrule
\end{tabular}
\caption{Summary of the repository's 7 real ADRs (\texttt{docs/adr/}).}
\label{tab:adrs}
\end{table}

\textbf{Corrección de integridad relevante:} ADR-007 afirmaba, en su redacción original, que las
imágenes base ya estaban ancladas por digest SHA-256 en \texttt{docker-compose.yml}. Verificado contra el
archivo real durante la Fase 8 de este proyecto: esa afirmación era falsa (las imágenes usaban solo tags
mutables). Se implementó el anclaje real en esa misma revisión, con los tres digests verificados vía
\texttt{docker pull}, en vez de simplemente corregir el texto de la ADR sin corregir el código que
describía. (Nota: ADR-006 y ADR-007 se renumeraron el 2026-09-01 --- antes eran ADR-003 y ADR-006
respectivamente --- para que coincidan con los temas específicos que exige la guía de la Entrega Final
en esos dos números; el contenido no cambió, solo el orden.)

\subsection{Modelo de datos}

El esquema relacional (\texttt{presus}, PostgreSQL 15) modela 13 entidades principales (usuarios,
solicitudes, anteproyectos, jurados, cronogramas, evaluaciones, actas, rúbricas, salas, notificaciones,
entre otras) con claves foráneas que garantizan integridad referencial. Los identificadores son
\texttt{BIGINT} autoincrementales (migrados desde UUID en una corrección temprana documentada en
OBS-01), y cuatro procedimientos almacenados PL/pgSQL puros (\texttt{sp\_calcular\_promedio\_evaluacion},
\texttt{sp\_asignar\_jurado\_masivo}, \texttt{sp\_firmar\_acta\_digital},
\texttt{sp\_generar\_reporte\_defensas}) encapsulan la lógica académica que involucra uniones
multi-tabla, versionados vía Flyway en \texttt{backend/src/main/resources/db/migration/}.
\newpage
